% 03-protocol.tex -- Evaluation protocol. Source of every statement: 05-protocol/protocol.md (fixed 2026-11-10,
% before any run) and experiments/README.md (counts, verification). Numbers in this file are measured, not pending.
\section{Evaluation Protocol}\label{sec:protocol}

The protocol was written, reviewed and frozen before any detector was scored on a test split; later changes are
logged with their reason. Its question is: under one leak-controlled protocol, how well do detectors separate
each Sybil attack from benign traffic, as a function of what the training data is assumed to be (benign only,
an unlabeled mix, or labeled)?

\subsection{Data and Unit of Analysis}\label{sec:data}

We use the four Sybil attacks of VeReMi-Extension~\cite{kamel2020veremi} (GridSybil, DataReplaySybil,
DoSRandomSybil, DoSDisruptiveSybil), each simulated for two periods (0709 and 1416), i.e., eight scenario
folders. Beacons are sent at 1\,Hz. Every vehicle logs the messages it receives together with its own position
and velocity at reception.

\textbf{Link.} The unit is receiver-side: a \emph{link} is the time-ordered sequence of messages that one
receiving vehicle heard from one sender pseudonym. Every received copy is kept, so one broadcast heard by five
vehicles contributes to five links. This is what a detector deployed on a single vehicle observes; it needs no
identity beyond the pseudonym.

\textbf{Window.} Each link is cut into consecutive, non-overlapping pieces of at most $T=32$ messages, so every
message is in exactly one window. A window is a $32\times13$ matrix, zero-padded at the end, with a prefix mask
that marks its $n\in\{1,\dots,32\}$ real rows. The 13 features per message are the position, velocity and
acceleration claimed by the sender ($x$, $y$ each), the receiver's own position and velocity ($x$, $y$ each), the
range and the bearing from the receiver to the claimed position, and the logarithm of the time since the previous
message on the link (missing for the first message). Time of day, the true sender identifier and all ground-truth
fields are never inputs. Features are standardized with the mean and standard deviation of the real rows of the
benign training windows; padding rows and the missing inter-arrival value are zero after standardization.

\textbf{Labels.} A window inherits the attack code of its sender vehicle's trace file. Benign windows come from
senders that are never attackers in their folder. Labels select the training pool of a regime
(Section~\ref{sec:regimes}) and are otherwise read only by the evaluation module, after all scores are written.

\subsection{Sender-Held-Out Split}\label{sec:split}

The four attack families reuse the same underlying traffic: the benign vehicles of one family reappear, with the
same trajectories, in the others. A split made inside one scenario folder, or a random split over windows,
therefore places near-copies of test trajectories in the training set, which is a train--test dependence in the
sense of~\cite{kapoor2023leakage,arp2022dos}. \todo{Optional supporting number from our earlier audit of a
within-scenario vehicle split: 89.7\% of benign test identities had a near-copy in train
(docs/research-notes/data\_understanding/cross-scenario-leak.md). Decide whether it goes here or in the
background section.}

We therefore split by \emph{vehicle key} (period, physical vehicle) across all eight folders: 80\,/\,10\,/\,10
for train\,/\,validation\,/\,test, stratified by period and class of the key, with one fixed seed. All windows
sent by a vehicle, in every folder and at every receiver, fall in one part. The split has 5{,}588\,/\,696\,/\,750
keys with windows, and no key occurs in two parts. The benign test keys and the attack test keys are disjoint in
each of the four tasks. The split is \emph{sender-held-out, not receiver-held-out}: the four receiver features
come from vehicles that also hear training senders. One split, one normalization and one test set are used for
every detector (Table~\ref{tab:data}).

\begin{table}[t]
\caption{Windows of $32\times13$ per split and class. ``1--3 rows'' is the share of windows with at most three
real rows. GridSybil train and validation windows are used only by the regimes that include attacks.}
\label{tab:data}
\centering
\footnotesize
\setlength{\tabcolsep}{4pt}
\begin{tabular}{@{}llrrr@{}}
\toprule
Part & Class & Windows & Vehicles & 1--3 rows \\
\midrule
Train & Benign        & 502{,}773 & 4{,}206 & 13.7\% \\
      & GridSybil     & 229{,}080 & 1{,}382 & 18.5\% \\
\midrule
Val.  & Benign        &  65{,}313 &   524 & 13.9\% \\
      & GridSybil     &  28{,}374 &   172 & 19.1\% \\
\midrule
Test  & Benign        &  60{,}142 &   525 & 13.9\% \\
      & GridSybil     &  28{,}087 &   172 & 20.4\% \\
      & DataReplay    &  76{,}283 &   173 & 94.9\% \\
      & DoSRandom     & 199{,}834 &   173 & 99.7\% \\
      & DoSDisruptive & 199{,}826 &   173 & 99.7\% \\
\bottomrule
\end{tabular}
\end{table}

The GridSybil training and validation windows were prepared with a different normalization; they are assigned
to the parts by vehicle key and re-standardized with the benign training statistics. As a check, the GridSybil
windows of the test vehicles derived this way equal the test split built directly: 28{,}087 windows, identical
identifiers and masks, maximum absolute difference $1.9\times10^{-4}$ in normalized units (rounding). All
285{,}541 GridSybil windows are assigned to exactly one part.

Table~\ref{tab:data} shows the property that shapes the evaluation. DataReplay and both DoS attacks change
pseudonym after one or two messages, so 94.9--99.7\% of their windows have at most three rows, against
13.9\% of benign test windows.

\subsection{Training Regimes}\label{sec:regimes}

A regime states what the training data is assumed to be (Table~\ref{tab:regimes}).
\emph{Benign-only} is the usual one-class assumption; selecting benign windows uses labels.
\emph{Mix-all} and \emph{mix-gridfolders} pool benign and GridSybil windows without labels.
\emph{Benign-gridfolders} is the control for mix-gridfolders: the same benign windows without the attack.
\emph{Supervised} uses the mix-all pool with labels; GridSybil is the only attack ever seen in training.

\begin{table}[t]
\caption{Training regimes. The validation pool of a regime has the same composition, from the validation
vehicles. ``Grid folders'' are the two GridSybil scenario folders.}
\label{tab:regimes}
\centering
\footnotesize
\setlength{\tabcolsep}{3pt}
\begin{tabular}{@{}p{1.9cm}p{2.55cm}rrr@{}}
\toprule
Regime & Training pool & Windows & \shortstack[r]{Grid-\\Sybil} & \shortstack[r]{Val.\\windows} \\
\midrule
Benign-only        & benign, 8 folders          & 502{,}773 & 0      & 65{,}313 \\
Benign-gridfolders & benign, grid folders       & 125{,}749 & 0      & 15{,}893 \\
Mix-all            & benign (8) + GridSybil     & 731{,}853 & 31.3\% & 93{,}687 \\
Mix-gridfolders    & benign + GridSybil, grid folders & 354{,}829 & 64.6\% & 44{,}267 \\
Supervised         & as mix-all, labeled        & 731{,}853 & 31.3\% & 93{,}687 \\
\bottomrule
\end{tabular}
\end{table}

The mix regimes are label-free end to end. Their inputs are re-standardized with the statistics of their own
training pool, model selection uses a validation pool of the same composition (30.3\% and 64.1\% GridSybil), and
their primary operating point is a percentile of that mixed validation pool. Pooled sets are stored in one fixed
random order, and detectors of these regimes receive features and mask only.

Mix-all and mix-gridfolders contain the same GridSybil windows and differ in the amount of benign traffic, which
itself repeats across families. They are two confounded points, not a dose--response curve. In the two GridSybil
folders the attack is the majority class (64.6\% of training windows), which is outside the small-contamination
setting usually studied~\cite{zhou2017anomaly,qiu2022latent}.

DataReplay and the DoS attacks are never in any training or validation pool. The supervised result on them is
\emph{transfer from GridSybil}, read within length strata; it supports no general claim about unseen attacks.

\subsection{Evaluation}\label{sec:evaluation}

\textbf{Tasks and metrics.} There are four binary tasks: the benign test windows against the test windows of
each attack. The attack shares are 0.318 (GridSybil), 0.559 (DataReplay) and 0.769 (each DoS attack). Every
detector outputs a score per window, higher meaning more anomalous. We report the area under the ROC curve
(AUROC), the average precision with the attack share stated, and the attack recall and benign false-positive rate
at fixed operating points.

\textbf{Operating points.} Thresholds come from validation scores only; an alarm is raised when the score exceeds
the threshold. The primary threshold is the 95th percentile of the benign validation scores for the benign and
supervised regimes, and the 95th percentile of the regime's own mixed validation pool for the mix regimes. For
the mix regimes the benign-validation threshold needs labels and is reported as an ``oracle calibration'' column
only. The 99th percentile of the benign validation scores is reported as a low false-positive point.

\textbf{Length control.} Every detector can see the window length (Section~\ref{sec:detectors}), and length alone
separates three of the four attacks (Table~\ref{tab:data}). Length is therefore controlled in the evaluation, in
three ways. First, a \emph{length-only} score $s(w)=-n(w)$, shorter meaning more anomalous, is evaluated like any
detector; its direction was fixed in the protocol and never chosen on labels. Second, the AUROC is reported
inside strata of the exact number of real rows, $n=1$, $n=2$, $n=3$, then $4\le n\le16$ and $n\ge17$ (and the
coarse strata $n\le3$, $n\ge4$). Third, we report a \emph{length-matched AUROC}. Let $B_n$ and $A_n$ be the
numbers of benign and attack test windows with exactly $n$ real rows and $\mathrm{AUROC}_n(s)$ the AUROC of score
$s$ on those windows. Then
\begin{equation}\label{eq:lm}
\mathrm{AUROC}_{\mathrm{LM}}(s)=
\frac{\sum_{n=1}^{32} B_n A_n\,\mathrm{AUROC}_n(s)}{\sum_{n=1}^{32} B_n A_n},
\end{equation}
where lengths with $B_nA_n=0$ contribute nothing. It is the probability that an attack window scores higher than
a benign window \emph{of the same length} (ties count one half), over all same-length pairs. It is 0.5 for the
length-only score by construction. A stratum with fewer than 200 attack windows or fewer than 20 attack vehicles
is not reported. A variant weights each length by $A_n$ alone, the number of attack windows of that length,
instead of $B_nA_n$; it is also 0.5 for the length-only score. The preliminary numbers of
Section~\ref{sec:res-prelim} use this attack-weighted variant, because that is what the earlier studies computed.

The length-matched AUROC is a methodological point of this paper, not only a control. A score can carry
window length without being a length feature: a reconstruction error, for example, is larger when a window has
few rows, because there is less context to reconstruct from. On the pseudonym-rotating attacks, where the attack
windows are short, such a score reaches a pooled AUROC well above 0.5 without separating attack from benign
windows of the same length. A detector can also score below the length-only control and still carry length;
comparing it with the control alone does not reveal this. Matching on length removes both effects
(Section~\ref{sec:disc-length} gives two examples from the earlier studies).

\textbf{Headline per attack.} For GridSybil the headline is the pooled AUROC, with the length-matched AUROC next
to it. For DataReplay and the DoS attacks a short pseudonym life is itself the observable signature of the
attack, and a pooled AUROC mostly measures length. Their headline is the length-matched AUROC and the paired
difference to length-only; the pooled AUROC is secondary.

\textbf{Seeds and uncertainty.} Stochastic detectors are run with seeds 0--4; we report the mean and the sample
standard deviation over seeds, which with five seeds is a rough estimate only. Deterministic detectors are run
once and marked. A 95\% percentile interval is obtained from a cluster bootstrap over test vehicle keys: benign
keys and attack keys are resampled separately, with 1{,}000 resamples shared by all detectors; in each resample
the metric is computed per seed and averaged. Resampling vehicle keys keeps all windows of a vehicle together,
including the copies of its traffic in the other families. The interval covers test sampling only: there is one
split, and thresholds are held fixed.

\textbf{Paired comparisons.} Because the resamples are shared, differences are paired. Five primary comparisons
were fixed in advance: (C1)~each detector minus length-only, per attack; (C2)~benign-only minus mix-all, same
detector, GridSybil; (C3)~benign-gridfolders minus mix-gridfolders, same detector, GridSybil; (C4)~the supervised
autoencoder with classifier head minus the best benign-only detector, GridSybil; (C5)~the convolutional
autoencoder minus each reference detector, benign-only, GridSybil.

\textbf{Sensitivity rows.} Three subsets are re-evaluated from the same score files, without retraining: the
1416 period only; the benign test windows of the GridSybil 0709 folder excluded; and the benign test set
restricted to the attack's own family folders. The second addresses a known limitation of the split
(Section~\ref{sec:limitations}).

\subsection{Records per Run}\label{sec:records}

Each run (regime, detector, seed) writes its configuration and hyperparameters, its environment (package
versions, device, git commit and dirty flag, a fingerprint of the data cache, the SHA-256 of every source file
used), the training history, one score file per validation and test split aligned with the window identifiers,
and a summary without any metric. All tables are rebuilt from the score files by one evaluation module, the only
code that knows which split is which class.
